\documentclass[10pt,a4paper]{article}
\usepackage[utf8]{inputenc}
\usepackage[T1]{fontenc}
\usepackage[english]{babel}
\usepackage{geometry}
\geometry{margin=1.8cm}
\usepackage{enumitem}
\usepackage{booktabs}
\usepackage{fancyhdr}
\usepackage{lastpage}
\usepackage{tcolorbox}

\pagestyle{fancy}
\fancyhf{}
\rfoot{\thepage/\pageref{LastPage}}
\lhead{\small TOEFL ITP 2026 --- Day 6: Mini-Test}
\rhead{\small MPrometeo}
\renewcommand{\headrulewidth}{0.4pt}
\setlength{\parindent}{0pt}
\setlength{\parskip}{5pt}

\begin{document}

\begin{center}
{\LARGE \textbf{DAY 6 --- Wednesday Mar 4}}\\[4pt]
{\large Mock Mini-Test}\\[2pt]
{\normalsize Theme: Pynchon's Dark Humor \& Absurdism}\\[2pt]
\rule{\textwidth}{0.8pt}
\end{center}

\section*{Context: Pynchon's Humor}

\begin{tcolorbox}[colback=white, colframe=black]
\small
Pynchon is one of the few ``serious'' novelists who is genuinely
funny. His humor operates on multiple levels simultaneously: slapstick
physical comedy, sexual absurdity, dark political satire, and
linguistic wordplay. In \emph{Gravity's Rainbow}, a character named
Pirate Prentice makes a giant breakfast of bananas on a London
rooftop during the Blitz. In \emph{Vineland}, a character called
the Thanatoids are people who are technically dead but too lazy to
lie down. Characters have names like Oedipa Maas, Tyrone Slothrop,
and Mucho Maas. There is a recurring song about a man who is in
love with his own nose.

This humor is not decorative. It is structural. Pynchon uses
absurdity to make unbearable subjects --- war, death, state violence,
the meaninglessness of modern life --- bearable enough to examine.
The laughter is the spoonful of sugar that helps the existential
dread go down.
\end{tcolorbox}

% ============================================================
\section*{Listening Mini-Test (10 questions)}
% ============================================================

{\small\textit{Use TTS. Don't read while listening.}}

\begin{enumerate}[label=\textbf{\arabic*.}, leftmargin=*]

\item M: I tried to install Linux on my laptop last night.\\
W: Let me guess --- you spent four hours configuring Wi-Fi drivers.\\
Q: What does the woman imply?
\begin{enumerate}[label=(\Alph*)]
\item Linux is easy to install
\item Linux often has driver compatibility issues
\item She doesn't know about computers
\item The man succeeded quickly
\end{enumerate}

\item W: The professor said the final is ``comprehensive.''\\
M: Great. So basically everything we've ever learned. No pressure.\\
Q: The man's tone suggests:
\begin{enumerate}[label=(\Alph*)]
\item He's excited about the exam
\item He's being sarcastic about the difficulty
\item He thinks it will be easy
\item He already studied everything
\end{enumerate}

\item M: Did you read Pynchon's description of the banana breakfast
in \emph{Gravity's Rainbow}?\\
W: Who writes three pages about bananas during a bombing and makes
it the funniest thing you've ever read?\\
Q: What does the woman think about the passage?
\begin{enumerate}[label=(\Alph*)]
\item It was boring and irrelevant
\item It was absurdly funny despite its dark setting
\item She didn't understand it
\item Bananas aren't funny
\end{enumerate}

\item W: My code compiled on the first try today.\\
M: Quick, buy a lottery ticket. That kind of luck doesn't come
around often.\\
Q: What does the man imply?
\begin{enumerate}[label=(\Alph*)]
\item She should gamble more
\item Code compiling on the first try is extremely rare
\item He doesn't believe her
\item He wants to buy a ticket
\end{enumerate}

\item M: The university is replacing all Windows computers with
Chromebooks.\\
W: So we're going from an operating system that crashes to one that
barely qualifies as an operating system.\\
Q: What does the woman think?
\begin{enumerate}[label=(\Alph*)]
\item Both options are inadequate
\item Chromebooks are superior
\item Windows never crashes
\item She prefers Chromebooks
\end{enumerate}

\item W: I told my advisor I wanted to write my thesis on Pynchon.\\
M: What did she say?\\
W: She asked if I enjoyed suffering.\\
Q: What does the advisor's response suggest?
\begin{enumerate}[label=(\Alph*)]
\item She hates Pynchon personally
\item Writing about Pynchon is extremely challenging
\item Pynchon is easy to analyze
\item The thesis was approved immediately
\end{enumerate}

\item M: How's the new Maple license working out?\\
W: Put it this way --- I've spent more time figuring out the
license manager than actually doing math.\\
Q: What is the woman's complaint?
\begin{enumerate}[label=(\Alph*)]
\item Maple is too expensive
\item The licensing system is more complicated than the software itself
\item She prefers Mathematica
\item The math is too difficult
\end{enumerate}

\item W: Did you understand the lecture on entropy?\\
M: I understood the word ``the'' and everything went downhill from
there.\\
Q: What does the man mean?
\begin{enumerate}[label=(\Alph*)]
\item He understood most of the lecture
\item He understood almost nothing
\item The lecture was about skiing
\item He's an expert on entropy
\end{enumerate}

\item M: Pynchon hasn't given an interview in over 60 years.\\
W: At this point, his reclusiveness is its own art form.\\
Q: What does the woman mean?
\begin{enumerate}[label=(\Alph*)]
\item Pynchon should do more interviews
\item His avoidance of publicity has become remarkable in itself
\item She thinks he's rude
\item He's not actually reclusive
\end{enumerate}

\item W: I switched from Mathematica to open-source alternatives.\\
M: How's that going?\\
W: I have more freedom and less hair.\\
Q: What does the woman imply?
\begin{enumerate}[label=(\Alph*)]
\item She got a haircut
\item Open-source is easy and relaxing
\item The switch gave her freedom but caused frustration
\item She regrets nothing
\end{enumerate}

\end{enumerate}

% ============================================================
\section*{Structure Mini-Test (10 questions)}
% ============================================================

\begin{enumerate}[label=\textbf{\arabic*.}, start=11, leftmargin=*]

\item Pynchon, \rule{2cm}{0.4pt} novels are famously difficult,
is also one of America's funniest writers.
\begin{enumerate}[label=(\Alph*)]
\item who \item whose \item which \item whom
\end{enumerate}

\item The humor in \emph{Gravity's Rainbow} makes the horror
\rule{2cm}{0.4pt} to process emotionally.
\begin{enumerate}[label=(\Alph*)]
\item enough bearable \item bearable enough
\item more bearable enough \item bearable than
\end{enumerate}

\item Had Pynchon \rule{2cm}{0.4pt} a public figure, his mystique
would not exist.
\begin{enumerate}[label=(\Alph*)]
\item being \item been \item be \item was
\end{enumerate}

\item \textit{Find error:}
(A)\textbf{Despite} of his (B)\textbf{reclusiveness}, Pynchon
(C)\textbf{remains} one of the most (D)\textbf{discussed} American
authors.

\item It is essential that dark humor \rule{2cm}{0.4pt} understood
in its cultural context.
\begin{enumerate}[label=(\Alph*)]
\item is \item be \item will be \item being
\end{enumerate}

\item \textit{Find error:}
Pynchon's characters (A)\textbf{has} names that (B)\textbf{sound}
(C)\textbf{deliberately} (D)\textbf{absurd}.

\item The Thanatoids, \rule{2cm}{0.4pt} are technically dead,
represent Pynchon's darkest joke about American passivity.
\begin{enumerate}[label=(\Alph*)]
\item which \item who \item whose \item whom
\end{enumerate}

\item Seldom \rule{2cm}{0.4pt} a novelist combined intellectual
depth with such crude humor so effectively.
\begin{enumerate}[label=(\Alph*)]
\item a novelist has \item has a novelist
\item a novelist had \item did a novelist
\end{enumerate}

\item \textit{Find error:}
Not only (A)\textbf{Pynchon is} a (B)\textbf{brilliant} stylist,
but he is also (C)\textbf{genuinely} (D)\textbf{hilarious}.

\item The professor recommended that every student \rule{2cm}{0.4pt}
at least one Pynchon novel before graduating.
\begin{enumerate}[label=(\Alph*)]
\item reads \item read \item will read \item reading
\end{enumerate}
\end{enumerate}

% ============================================================
\section*{Reading Mini-Test: Pynchon's Absurd Humor ($\sim$200 words)}
% ============================================================

\begin{tcolorbox}[colback=white, colframe=black]
\small
There is a particular kind of comedy that only works in proximity to
disaster. The banana breakfast scene in \emph{Gravity's Rainbow} is
a perfect example: Pirate Prentice prepares an elaborate meal of
bananas --- fried, broiled, flambeed --- on a London rooftop while
German V-2 rockets fall on the city below. The juxtaposition is
absurd. It is also, somehow, deeply human.

Pynchon understood that humor and horror are not opposites. They are
siblings --- born from the same recognition that the universe
operates according to rules we cannot fully understand or control.
When a character in \emph{Vineland} is described as being ``too dead
to lie down,'' the joke works precisely because death is not funny.
The laughter comes from the gap between what death \emph{should}
mean and how casually Pynchon's characters treat it.

This technique has literary ancestors: Rabelais, Swift, Voltaire,
Kafka. But Pynchon adds a distinctly American flavor --- the humor
of a country that invented both the atomic bomb and the whoopee
cushion. It is comedy that does not deny suffering but refuses to
let suffering have the last word.
\end{tcolorbox}

\begin{enumerate}[label=\textbf{\arabic*.}, start=21]
\item The banana breakfast scene works because:
\begin{enumerate}[label=(\Alph*)]
\item Bananas are inherently funny
\item The absurd normality contrasts with the surrounding danger
\item The rockets stopped falling
\item Prentice is a professional chef
\end{enumerate}

\item ``Juxtaposition'' is closest in meaning to:
\begin{enumerate}[label=(\Alph*)]
\item separation \item placing contrasting things side by side
\item repetition \item explanation
\end{enumerate}

\item According to the passage, humor and horror are:
\begin{enumerate}[label=(\Alph*)]
\item Complete opposites
\item Related responses to an uncontrollable universe
\item Irrelevant to literature
\item Only found in American fiction
\end{enumerate}

\item The ``American flavor'' Pynchon adds refers to:
\begin{enumerate}[label=(\Alph*)]
\item Using only American settings
\item The irony of a nation capable of both destruction and silliness
\item Writing only about food
\item Avoiding European influences
\end{enumerate}

\item The final sentence means:
\begin{enumerate}[label=(\Alph*)]
\item Comedy replaces suffering entirely
\item Comedy acknowledges suffering but refuses to surrender to it
\item Suffering is not real
\item Comedy has no purpose
\end{enumerate}
\end{enumerate}

\textbf{All Answers:}

\textbf{Listening:} 1-B, 2-B, 3-B, 4-B, 5-A, 6-B, 7-B, 8-B, 9-B, 10-C

\textbf{Structure:} 11-B, 12-B, 13-B, 14-A (Despite OF $\to$ Despite),
15-B (subjunctive), 16-A (has $\to$ have), 17-B (who for people),
18-B (inversion: has a novelist), 19-A (Not only IS Pynchon),
20-B (subjunctive: read)

\textbf{Reading:} 21-B, 22-B, 23-B, 24-B, 25-B

\end{document}